% Options for packages loaded elsewhere
\PassOptionsToPackage{unicode}{hyperref}
\PassOptionsToPackage{hyphens}{url}
%
\documentclass[
  11pt,
]{article}
\usepackage{amsmath,amssymb}
\usepackage{iftex}
\ifPDFTeX
  \usepackage[T1]{fontenc}
  \usepackage[utf8]{inputenc}
  \usepackage{textcomp} % provide euro and other symbols
\else % if luatex or xetex
  \usepackage{unicode-math} % this also loads fontspec
  \defaultfontfeatures{Scale=MatchLowercase}
  \defaultfontfeatures[\rmfamily]{Ligatures=TeX,Scale=1}
\fi
\usepackage{lmodern}
\ifPDFTeX\else
  % xetex/luatex font selection
\fi
% Use upquote if available, for straight quotes in verbatim environments
\IfFileExists{upquote.sty}{\usepackage{upquote}}{}
\IfFileExists{microtype.sty}{% use microtype if available
  \usepackage[]{microtype}
  \UseMicrotypeSet[protrusion]{basicmath} % disable protrusion for tt fonts
}{}
\makeatletter
\@ifundefined{KOMAClassName}{% if non-KOMA class
  \IfFileExists{parskip.sty}{%
    \usepackage{parskip}
  }{% else
    \setlength{\parindent}{0pt}
    \setlength{\parskip}{6pt plus 2pt minus 1pt}}
}{% if KOMA class
  \KOMAoptions{parskip=half}}
\makeatother
\usepackage{xcolor}
\usepackage[margin=0.9in]{geometry}
\setlength{\emergencystretch}{3em} % prevent overfull lines
\providecommand{\tightlist}{%
  \setlength{\itemsep}{0pt}\setlength{\parskip}{0pt}}
\setcounter{secnumdepth}{-\maxdimen} % remove section numbering
\ifLuaTeX
  \usepackage{selnolig}  % disable illegal ligatures
\fi
\usepackage{bookmark}
\IfFileExists{xurl.sty}{\usepackage{xurl}}{} % add URL line breaks if available
\urlstyle{same}
\hypersetup{
  pdftitle={Pair-block closure in the one-isolate case},
  pdfauthor={Aafi Khan Malek},
  hidelinks,
  pdfcreator={LaTeX via pandoc}}

\title{Pair-block closure in the one-isolate case}
\author{Aafi Khan Malek\\[-2pt]\small Independent researcher, Ahmedabad, India\\[-2pt]\small\href{mailto:aafimalek2032@gmail.com}{aafimalek2032@gmail.com}}
\date{27 November 2026}

\begin{document}
\maketitle

\textbf{Scope.} Consider a ground-minimal weak counterexample at height
at most six, with exactly one isolated point and at least one forcing
edge, in the terminology of
\href{https://github.com/Aafimalek/novel-discoveries/blob/research/parallel-foundations/frankl-height-six/further-reductions.md}{the
project}. The argument below gives at least \textbf{14 forcing stars}.
The incidence argument through ten is now checked end to end in Lean.
The exclusions of eleven, twelve, and thirteen below are a written
proof, not yet checked end to end in Lean. No case with a different
number of isolated points is covered.

\subsection{Reduction to a design and a fiber
bound}\label{reduction-to-a-design-and-a-fiber-bound}

Let \(K\) be the set of star centers, \(k=|K|\), and \(z\) the isolated
point. The checked one-isolate result gives \(k\ge 9\). Each star has
one leaf: otherwise the checked multi-leaf theorem, using a different
center, would give at least four isolates. Thus the ground has
\(n=2k+1\) points.

For distinct \(p,q\in K\), let \(T_{pq}\) be the largest family member
avoiding the two corresponding leaves. In this fiber, every star is
present or absent as a unit. Let \(B_{pq}\) be the set of centers whose
stars are absent from \(T_{pq}\).

The checked pair-exchange lemma gives \(T_{pq}=T_{pr}\) whenever
\(r\in B_{pq}\setminus\{p\}\). For any distinct \(r,s\in B_{pq}\), first
exchange at \(p\) to use \(r\), then at \(r\) to use \(s\) (if either is
\(p\), one exchange suffices). Therefore \(T_{rs}=T_{pq}\). The distinct
sets \(B_{pq}\) form a collection \(\mathcal{B}\) of blocks: each pair
appears in exactly one block, and distinct blocks meet in at most one
center. Since equal blocks contain a common pair, different maxima have
different blocks.

Let \(W(p,q)\) count blocks containing neither \(p\) nor \(q\), and put
\(H(m)=1+m+\binom m2\). The following two-sided bound will be used:

\[
 3k+W(p,q) \le |\mathcal{F}_{pq}| \le H(k+1-|B_{pq}|),
 \tag{1}
\]

where \(\mathcal{F}_{pq}\) is the family avoiding the two leaves. For
the upper bound, compress each surviving star to one coordinate,
retaining the isolate if present; the compressed family is union-closed
of height at most three. Its support is of size at most
\(k+1-|B_{pq}|\). For the lower bound, the full ground and all one-point
avoidance maxima except those for the two prescribed leaves supply
\(n-1=2k\) distinct both-containing members. A further \(k\) are the
distinct complements of a center and \(z\). These are noncanonical by
the checked source--isolate witness lemma. Finally each block disjoint
from \(\{p,q\}\) gives its own both-containing maximum, omitting at
least four points and hence distinct from those earlier members. The
one-point maxima omit at most two points by the star structure. Each
leaf has frequency at most half, so there are at least as many
neither-containing as both-containing members.

\subsection{Two elementary bounds}\label{two-elementary-bounds}

\textbf{Incidences.} Suppose all blocks have size at most \(r\). For an
outside vertex \(v\in K\setminus\{p,q\}\), its block through \(p\) and
its block through \(q\) can account for at most \(2(r-2)\) other outside
vertices altogether. Hence at least \(k-3-2(r-2)\) of its outside
partners use blocks disjoint from \(\{p,q\}\). A counted block through
\(v\) has at most \(r-1\) such partners. Consequently, for every
positive integer \(d\) such that

\[
   (r-1)(d-1)<k-3-2(r-2),
\]

double-counting incidences with the \(k-2\) outside vertices gives
\(rW(p,q)\ge (k-2)d\).

\textbf{Closure.} Fix a block \(B\), and form a simple graph on
\(K\setminus B\), with edge \(uv\) exactly when the block of \(u,v\)
meets \(B\). Any member \(A\) avoiding a pair from \(B\) omits all stars
of \(B\). Write \(L(A)\) for the stars it omits. In general if
\(r,s\in L(A)\) are distinct, then \(A\subseteq T_{rs}\) and
\(B_{rs}\subseteq L(A)\). In particular, if \(uv\) is a graph edge and
\(u\in L(A)\), use the pair consisting of \(u\) and the unique point of
\(B\) in \(B_{uv}\) to conclude \(v\in L(A)\). Thus \(L(A)\setminus B\)
is a union of connected components. By twinning, \(A\) is determined by
these components and the isolate bit. If the graph has \(c\) components,
\(|\mathcal{F}_{pq}|\le 2^{c+1}\) for any pair in \(B\). Whenever we use
this closure bound below, \(3k>32\). It suffices to find a block with
\(c\le 4\).

For a fixed \(b\in B\), its other incident blocks partition the
\(h=k-|B|\) outsiders into groups, with \(g_b\) groups. Those groups
yield \(E_b=\sum_S\binom{|S|}{2}\) link edges, and \(g_b\ge h-E_b\).
Each graph edge is counted at just one \(b\); write \(E=\sum_b E_b\).
For distinct \(p,q\in B\), the groups at each remaining \(b\) give
\(g_b\) distinct blocks that avoid \(p,q\). The other such blocks are
disjoint from \(B\); they cover the \(\binom h2-E\) outside pairs not in
the graph. When all blocks have size at most \(r\), this proves

\[
 W(p,q)\ge \sum_{b\in B\setminus\{p,q\}}(h-E_b)
           +\left\lceil\frac{\binom h2-E}{\binom r2}\right\rceil .
 \tag{2}
\]

For graphs on eight and nine vertices, respectively, five or more
components imply at most six and ten edges. Indeed, the total excess of
component sizes over one is at most three or four, and concentrating it
in one component maximizes the sum of the clique edge bounds.

\subsection{Excluding nine and ten}\label{excluding-nine-and-ten}

By (1) with \(W\ge 0\), every block has size at most three. The
incidence bound with \(r=3,d=2\) at nine, and \(r=3,d=3\) at ten, yields
\(W\ge 5,8\) respectively. This raises the lower bound to \(32,38\);
those numbers exceed \(H(k-2)=29,37\) respectively. Therefore all blocks
are pairs (and the isolate must survive in every pair maximum). Now
\(W=\binom{k-2}{2}\), and the bounds (1) are inconsistent: \(48>37\) at
nine, \(58>46\) at ten.

\subsection{Reducing eleven, twelve and thirteen to blocks of size at
most
three}\label{reducing-eleven-twelve-and-thirteen-to-blocks-of-size-at-most-three}

At eleven, (1) with \(W\ge 0\) leaves size at most four. With
\(r=4,d=2\) the incidence bound gives \(W\ge 5\), so
\(3k+W\ge 38>H(8)=37\), excluding size four.

At twelve, (1) with \(W\ge 0\) leaves size at most five. With
\(r=5,d=1\) the incidence bound gives \(W\ge 2\), so
\(3k+W\ge 38>H(8)=37\), excluding size five. If a block of size four
remains, it has eight outsiders. If \(E\ge 7\), closure gives the
contradiction. If \(E\le 6\), for any pair in the block (2) gives
\(W\ge 16-6+\lceil 22/6\rceil=14\), but (1) gives \(W\le H(9)-36=10\).

At thirteen, (1) with \(W\ge 0\) leaves size at most five. A block of
size five has eight outsiders. If \(E\ge 7\), use closure; otherwise (2)
gives \(W\ge 24-6+\lceil 22/10\rceil=21\), but (1) gives
\(W\le H(9)-39=7\). A remaining size-four block has nine outsiders. If
\(E\ge 11\), use closure. Otherwise the two centers \(b,c\) with the
smallest link-edge counts satisfy \(E_b+E_c\le 5\). Use the
complementary pair \(p,q\) in (2) with \(r=4\):
\(W\ge 18-5+\lceil 26/6\rceil=18\), but (1) gives \(W\le H(10)-39=17\).

\subsection{Finishing the three triple-block
cases}\label{finishing-the-three-triple-block-cases}

Suppose now \(k\in\{11,12,13\}\) and all blocks are pairs or triples. If
all blocks are pairs, \(W=\binom{k-2}{2}\) gives lower bounds
\(69,81,94\), exceeding the respective \(H(k-1)=56,67,79\). Let \(t>0\)
be the number of triple blocks and \(d_v\) the number using \(v\). Then
there are \(\binom k2-2t\) blocks in total and \(k-1-d_v\) through
\(v\). Inclusion--exclusion gives

\[
 W(p,q)=\binom{k-2}{2}-2t+d_p+d_q.
\]

For every pair in a triple, (1) gives \(W(p,q)\le U:=H(k-2)-3k\).
Summing over the \(3t\) different pairs in triples gives

\[
 2\sum_v d_v^2 \le 3t\left(2t+U-\binom{k-2}{2}\right).
\]

By Cauchy--Schwarz and \(\sum_v d_v=3t\), the left side is at least
\(18t^2/k\). This forces \(t\ge 16,17,18\) at eleven, twelve and
thirteen, respectively. Also
\(\sum_v d_v^2=\sum_{B:\,|B|=3}\sum_{b\in B}d_b\), so some triple block
has degree-sum at least \(\lceil 9t/k\rceil\). Its link graph has at
least \(\lceil 9t/k\rceil-3=11,10,10\) edges, respectively, from the
triples other than \(B\) meeting it. Edges from one fixed \(b\) are a
matching, so maximum degree is at most three.

For eight outsiders, eleven edges force at most four components. For
nine outsiders, if there are at least five components, their total
excess over one is at most four. For ten outsiders, it is at most five.
A degree-at-most-three component of size \(s=1,2,\dots,6\) has edge
bound \(0,1,3,6,7,9\), respectively, using the minimum of the clique and
degree-sum bounds. With excess at most four or five, these bounds sum to
at most seven or nine, respectively: each is at most \(2(s-1)\), and
equality at positive excess requires \(s=4\), so equality at total
excess four or five is impossible. In both cases ten edges force at most
four components. Closure now gives at most 32 pair-avoiding members,
contradicting \(3k\ge 33\). This finishes the written exclusion through
thirteen.

\subsection{Reproducibility}\label{reproducibility}

The checked structural hypotheses and transcript are in the project's
Lean directory. The formal lower bound of eleven is
\texttt{one\_isolate\_nonzero\_sources\_ge\_eleven} in
\texttt{FranklBlockWitnesses.lean}. The companion script checks the
numeric steps and the graph component edge bounds for the written
extension to fourteen. It is not a replacement for either the structural
proof or Lean.

ChatGPT assisted in proof development, drafting, and independent
checking. No Codex coding agent was used.

\end{document}
